\section{Barriers}\label{sec:barriers}

This section explains why the Tape Coding Conjecture is hard to prove and why a proof must use the
shape of tape networks. Its results are not formalized; each gets a complete proof.

\subsection{Multiplication}

By Harvey and van der Hoeven~\cite[Cor.~6.9]{HvdH2025}, a multiplication machine running in time
$M(m)$ yields a transposition machine for $n\times n$ matrices with $b$-bit entries running in time
$O(\ell_0\,\overline M(m))$, where $m\ge n^2b$, $\ell_0=\max(1,\lg^*n-\lg^*b)$ and
$\overline M(m)=m\max_{k\le m}M(k)/k$. For $b=\lceil\log_2n\rceil$ we have $\ell_0=1$, so a
linear-time multiplication machine gives a linear-time transposer. Hence $\neg$(T), and in particular
the Tape Coding Conjecture, implies that integer multiplication cannot be done in linear time on
multitape Turing machines. No superlinear lower bound for off-line multiplication on multitape
machines is known; the classical $\Omega(n\log n)$ bound of Paterson, Fischer and Meyer is for
on-line multiplication~\cite{PFM1974}.

\subsection{General graphs: the circuit barrier}

We use the following elementary model. Let $H$ be an undirected graph with source vertices
$a_0,\dots,a_{n-1}$ and sink vertices $b_0,\dots,b_{n-1}$, every edge of unit capacity. The
\emph{complete bipartite demand} asks to deliver, for every pair $(i,j)$, a message from $a_i$ to $b_j$.
\begin{itemize}[leftmargin=2em]
\item A \emph{routing at rate $\lambda$} is a fractional multicommodity flow that sends $\lambda$ units
from $a_i$ to $b_j$ for every pair, with total flow at most $1$ on every edge.
\item A \emph{causal code with $N$ rounds} fixes an acyclic orientation of $H$ and, in each round, sends
one bit along every edge in its orientation, each bit being a function (which may depend on the round) of the
bits that entered the tail in the same round and, at a source $a_i$, of $a_i$'s messages. Each pair $(i,j)$ is assigned a
round in which $b_j$ must output the message bit $m_{ij}$, a function of the bits entering $b_j$ in
that round; the messages $m_{ij}$ are arbitrary bits, and the code must work for all of them. Its
\emph{rate} is $1/N$ per pair and round.
\end{itemize}
In \cref{prop:rotation-barrier} we use three extensions of this model. A source may be attached at a set of
vertices: at each of them the code may use its messages, also in the output functions, a routing may start at any
of them, and $\dist_H(a_i,b_j)$ is the distance from the nearest. The functions may depend on the round. And edges
may carry symbols of $\kappa$ bits per round, against routings with capacity $\kappa$ per edge; then
\eqref{eq:weak-duality-graph} below reads $\lambda\sum_{i,j}\dist_H(a_i,b_j)\le\kappa\abs{E(H)}$.
The ratio of the rate of a causal code to the largest routing rate is its \emph{coding gain}. A routing
needs one unit of capacity per edge traversed, so for every routing,
\begin{equation}\label{eq:weak-duality-graph}
  \lambda\sum_{i,j}\dist_H(a_i,b_j)\;\le\;\abs{E(H)} .
\end{equation}

\begin{proposition}[Shift circuits are coding gains]\label{prop:circuit-barrier}
\notformalized
Let $C$ be a Boolean circuit with $s$ gates of fan-in at most two (arbitrary gate functions) whose
inputs are $x\in\{0,1\}^n$ and a shift input, such that for every $k<n$ some value of the shift input
makes the outputs equal the cyclic shift $x_{i-k\bmod n}$ at output $i$. Then there is a graph $H$ of
maximum degree at most $4$ with $O(s+n)$ vertices and a causal code with $n$ rounds for the complete
bipartite demand on $H$ whose coding gain is at least $c\,n\log n/(s+n)$, for an absolute constant
$c>0$ and all large $n$.
\end{proposition}

For example, a circuit multiplying two $n$-bit integers, followed by $n$ OR gates, is such a circuit:
for $y=2^k$ the $2n$-bit product $p=x\,2^k$ has $p_j$ and $p_{j+n}$ never both equal to $1$, so the
cyclic shift is $p_j\vee p_{j+n}$.

\begin{informalproof}
\emph{The graph.} Take the circuit as a directed acyclic graph on its gates and on the inputs that feed some gate (at most
$n+2s$ vertices; parallel edges allowed). Replace
every vertex of fan-out $f>2$ by a binary tree of $f-1$ copy vertices, so that in-degrees and
out-degrees are at most $2$; this adds at most $2s+n$ vertices. Let $H$ be the underlying undirected
graph, with $a_i$ the input vertex of $x_i$ and $b_j$ the vertex of output $j$. Its maximum degree is
at most $4$ and it has $O(s+n)$ vertices.

\emph{One code per shift.} Fix $k<n$ and fix the shift input to a value realizing the shift by $k$.
Every vertex then computes a fixed function of its in-edges, so the circuit, read in topological order
with one bit per edge, is a causal code on $H$ (with the circuit's orientation) that delivers $x_i$ from
$a_i$ to $b_{i+k\bmod n}$.

\emph{Time-sharing.} Run $n$ rounds; in round $k$ use the code for shift $k$ with the source $a_i$ sending
the message $m_{i,i+k\bmod n}$ as $x_i$. Every pair $(i,j)$ is served in round $j-i\bmod n$, and every edge
carries one bit per round. This is a causal code with $n$ rounds, of rate $1/n$.

\emph{The flow bound.} In a graph of maximum degree $4$ at most $4^{t+1}$ vertices lie within distance $t$
of a vertex. With $t=\lfloor\tfrac12\log_4n\rfloor-1$, at most $\sqrt n$ sinks lie within distance $t$ of
$a_i$, so
$\sum_{i,j}\dist(a_i,b_j)\ge n(n-\sqrt n)(t+1)\ge\tfrac15n^2\log_4n$ for large $n$. By
\eqref{eq:weak-duality-graph}, every routing has rate $\lambda\le5\abs{E(H)}/(n^2\log_4n)=O((s+n)/(n^2\log n))$.
The coding gain is therefore at least $(1/n)/\lambda=\Omega(n\log n/(s+n))$.
\end{informalproof}

So if, on every bounded-degree graph, complete bipartite exchange had coding gain at most $\gamma(n)$,
then every circuit computing all cyclic shifts, and every multiplication circuit, would have size
$\Omega(n\log n/\gamma(n))$. With $\gamma=o(\log n)$ this is a superlinear circuit lower bound for an
explicit function, which is far beyond current techniques: the best lower bounds for explicit functions
are $3.1n-o(n)$ over the full binary basis~\cite{LiYang2022}. A proof of the Tape Coding Conjecture must
therefore use the structure of tape networks: unless it also yields a superlinear circuit lower bound for an
explicit function, it cannot be a statement about bounded-degree networks in general. The single-shift version of this argument, conditional on the network coding conjecture, is due
to Afshani, Freksen, Kamma and Larsen~\cite{AFKL2019}, and the step from a small circuit to a coding gain
is also observed by Braverman and He~\cite{BravermanHe2026}.

Tape networks are not simple either. The chords of one work tape form two non-crossing families along the
time spine, excursions to the right and to the left, so a $k$-tape network has a book embedding with $2k$ pages,
while bounded-degree expanders have $3$-page book embeddings~\cite{DSW2016}. We expect, but do not prove, that after
spreading each vertex over a constant number of slots every page can be realized by the excursions of one work head,
so that tape networks with a constant number of work tapes contain subdivided bounded-degree expanders, on which cut
bounds would not exclude a code. This is the heuristic reason why we do not expect cut arguments alone to prove the
conjecture.

\subsection{Tape networks: an oblivious separation}

\emph{Programmable rotation} takes $x\in\{0,1\}^L$ and $\rho<L$ (in binary, $\clog L$ bits appended to $x$)
and outputs $x$ rotated by $\rho$. A non-oblivious machine computes it in time $O(L)$: copy $x$ to a work tape,
then count $\rho$ down in binary while moving a head $\rho$ cells (amortized $O(1)$ per decrement), and output the
work tape cyclically from that head. Oblivious machines can rotate in time $O(L\log L)$, by a barrel shifter or
by the oblivious simulation of Pippenger and Fischer~\cite{PippengerFischer1979}.

The \emph{bit-level network} of an oblivious machine $M$ at input length $L+\clog L$ has a vertex for every
step $\tau<T$; a spine edge $\tau\to\tau+1$; and, for every work tape and every pair of consecutive visits
$\sigma<\tau$ of its head to one cell, an edge $\sigma\to\tau$. The source of input bit $x_i$ is attached
at the steps where the input head reads cell $i$, and the sink of output bit $o$ is the step that appends
it. All of this is the same for every input, because $M$ is oblivious.

\begin{proposition}[Rotation codes on oblivious tape networks]\label{prop:rotation-barrier}
\notformalized
Let $M$ be an oblivious machine with $k$ work tapes that computes programmable rotation within $T(L)\le
L\log_2L$ steps. Then its bit-level network carries a causal code with $L$ rounds for the complete
bipartite demand between the input bits and the output bits whose coding gain over routing, with every edge
of capacity $\kappa=\max(\clog{(q+1)},\clog{(g+1)})$, is at least $c\,L\log L/T(L)$ for a constant $c>0$
depending only on $k$, $q$ and $g$, and all large $L$.
\end{proposition}

\begin{informalproof}
Fix $\rho<L$. With the bits of $\rho$ fixed, the run of $M$ is a causal code on its bit-level network: the spine
edge out of step $\tau$ carries the state, and the edge from one visit of a cell to the next carries the
symbol written at the first visit; every step computes its outgoing values from its incoming ones and the
input bit it reads. It delivers $x_i$ to output position $i-\rho\bmod L$. Time-sharing over the $L$ values of $\rho$, as in
\cref{prop:circuit-barrier}, serves every pair (input bit, output bit) once in $L$ rounds.

For the flow bound, every vertex has degree at most $D=2k+2$. Let $\mu_i$ be the number of steps at which
the input head reads cell $i$; then $\sum_i\mu_i\le T$. Within distance $t$ of the steps attaching $x_i$
lie at most $\mu_iD^{t+1}$ steps, hence at most $\mu_iD^{t+1}$ sinks. With
$t=\lfloor\tfrac12\log_DL\rfloor-1$, at most $T\sqrt L\le L^{3/2}\log_2L=o(L^2)$ pairs are within distance
$t$, so $\sum_{i,o}\dist\ge\tfrac12L^2(t+1)$ for large $L$. The network has at most $(k+1)T$ edges, each of
the stated capacity $\kappa$, so a routing has rate $\lambda\le(k+1)T\kappa/(\tfrac12L^2(t+1))=O(T/(L^2\log L))$,
while the code has rate $1/L$. The gain is $\Omega(L\log L/T)$.
\end{informalproof}

So a bit-level analogue of the Tape Coding Conjecture, asserting that the complete bipartite demand between
input and output bits on bit-level networks of oblivious machines running in linear time, at the capacities of
\cref{prop:rotation-barrier}, has coding gain $o(\log L)$, would imply that oblivious machines need time
$\omega(L)$ for programmable rotation, while non-oblivious machines need only $O(L)$. No superlinear
\emph{off-line} separation of oblivious and non-oblivious multitape machines is known: the optimality
claimed by Pippenger and Fischer is for on-line simulation~\cite{PippengerFischer1979}, and the known
superlinear lower bounds specific to oblivious machines are on-line or for one head~\cite{Vitanyi1984}. The block
version used for transposition and the bit version used here are not known to imply each other.
Independently of this calibration, by \cref{thm:oblivious} a proof of the conjecture would show that oblivious
machines cannot transpose in linear time, a lower bound that current techniques do not give. By contrast, Theorem~P and Theorem~M* are unconditional, and
the network coding conjecture, which would have given such bounds directly~\cite{AHJKL2006}, is
false~\cite{BravermanHe2026,ZhangLiLi2026}.
